% =============================================================================
\section{Ley de Gauss en dieléctricos: el vector desplazamiento}\label{sec:desplazamiento}
% =============================================================================

\subsection{Carga libre y carga ligada}\label{sec:gauss-D}

En presencia de dieléctricos, el campo electrostático lo crean dos tipos de carga: la libre, que
controlamos, y la ligada, que el material dispone por sí solo y que depende a su vez del campo. La ley de
Gauss del vacío sigue siendo válida si se cuentan ambas,
\begin{equation}\label{eq:gauss-total}
  \divg\vE=\frac{\rhof+\rhoP}{\eo},\qquad
  \oiint_S\vE\cdot\dS=\frac{\qfenc+\qPenc}{\eo},
\end{equation}
pero en la práctica no sirve de mucho: para usarla habría que conocer la carga de polarización, que es
precisamente lo que no sabemos hasta haber resuelto el problema. Conviene disponer de una ley que solo
involucre la carga libre.

\begin{lema}[Carga ligada encerrada]\label{lem:qP-enc}
Para cualquier superficie cerrada $S$, la carga de polarización que encierra es menos el flujo de $\vP$:
\begin{equation}\label{eq:qP-enc}
  \qPenc=-\oiint_S\vP\cdot\dS .
\end{equation}
\end{lema}

\begin{proof}
Sea $V_S$ la región encerrada por $S$ y $V$ la del dieléctrico. La carga ligada encerrada consta de la
volúmica de $V\cap V_S$ y de la superficial de la parte de $\partial V$ contenida en $V_S$. La frontera de
$V\cap V_S$ está formada por $\partial V\cap V_S$ y $S\cap V$; aplicando el teorema de la divergencia en esa región,
\[
  \int_{V\cap V_S}\rhoP\,\dd V=-\oiint_{\partial(V\cap V_S)}\vP\cdot\dS
  =-\int_{\partial V\cap V_S}\sigmaP\,\dd S-\int_{S\cap V}\vP\cdot\dS .
\]
Pasando el primer término de la derecha al otro miembro se obtiene $\qPenc=-\int_{S\cap V}\vP\cdot\dS$, que es
\eqref{eq:qP-enc} porque $\vP=\vcero$ fuera del dieléctrico.
\end{proof}

Llevando \eqref{eq:qP-enc}, o su forma local $\rhoP=-\divg\vP$, a \eqref{eq:gauss-total}, el flujo de $\vP$ pasa al
primer miembro:
\[
  \divg\vE=\frac{\rhof-\divg\vP}{\eo}
  \quad\Longrightarrow\quad
  \divg\big(\eo\vE+\vP\big)=\rhof .
\]
Esto sugiere agrupar los dos campos en uno solo.

\begin{definicion}[Vector desplazamiento eléctrico]\label{def:D}
Se llama \emph{vector desplazamiento eléctrico} (o inducción eléctrica) al campo
\begin{equation}\label{eq:D}
  \boxed{\;\vD(\vr)\defeq\eo\vE(\vr)+\vP(\vr)\; .}
\end{equation}
Sus unidades son las de $\vP$, $\si{C/m^2}$. El nombre se debe a Maxwell, que lo interpretaba como un
desplazamiento de las cargas del éter; hoy es simplemente una herramienta de cálculo.
\end{definicion}

\begin{teorema}[Ley de Gauss para $\vD$]\label{teo:gauss-D}
El vector desplazamiento satisface, en forma diferencial e integral,
\begin{equation}\label{eq:gauss-D}
  \divg\vD=\rhof,\qquad
  \oiint_S\vD\cdot\dS=\qfenc ,
\end{equation}
para cualquier superficie cerrada $S$: su flujo solo «ve» la carga libre.
\end{teorema}

\begin{proof}
La forma diferencial se acaba de deducir. Para la integral, por \eqref{eq:gauss-total} y el \cref{lem:qP-enc},
$\oiint_S\eo\vE\cdot\dS=\qfenc+\qPenc=\qfenc-\oiint_S\vP\cdot\dS$, que es \eqref{eq:gauss-D}. Ambas formas son equivalentes por el
teorema de la divergencia.
\end{proof}

\begin{observacion}[Lo que el teorema no dice]\label{obs:D-no-determinado}
La ley de Gauss para $\vD$ es tan útil como la del vacío \emph{cuando hay simetría suficiente}: plana, cilíndrica
o esférica. En general, conocer la divergencia de un campo no basta para determinarlo; hace falta también su
rotacional, y el de $\vD$ no es nulo (\cref{sec:rot-D}). En particular, que la carga libre encerrada sea nula
\emph{no} implica que $\vD$ sea nulo.
\end{observacion}

\subsection{Electretes}\label{sec:electretes}

\begin{definicion}[Electrete]\label{def:electrete}
Un \emph{electrete} es un material con una polarización permanente, que no se debe a un campo aplicado y no
desaparece al retirarlo. Los materiales ferroeléctricos son el ejemplo más importante; son el análogo eléctrico
de los imanes permanentes (ferromagnéticos) del magnetismo.
\end{definicion}

En un electrete $\vP$ es un dato del problema, y basta el \cref{teo:densidades} para calcular su campo. Los
problemas~\ref{prob:cilindro-polarizado}, \ref{prob:electrete-esferico}, \ref{prob:electrete-coaxial}
y~\ref{prob:electrete-condensador} son de este tipo.

\subsection{El rotacional de $\vD$}\label{sec:rot-D}

El campo electrostático es conservativo, $\rot\vE=\vcero$. Tomando el rotacional de \eqref{eq:D},
\begin{equation}\label{eq:rot-D}
  \rot\vD=\rot\vP .
\end{equation}
El desplazamiento no es, en general, conservativo: su circulación no se anula allí donde la tiene $\vP$.
Recordemos el teorema de Helmholtz~\cite{T0}.

\begin{lema}[Teorema de Helmholtz]\label{lem:helmholtz}
Un campo vectorial $\vF$ de clase $\mathcal C^1$ a trozos que tiende a cero en el infinito más deprisa que $1/r$
queda determinado de forma única por su divergencia y su rotacional (entendidos en sentido distribucional,
de modo que incluyen las discontinuidades de sus componentes normal y tangencial en las superficies).
\end{lema}

En una superficie en que $\vP$ salta de $\vP_1$ a $\vP_2$, la parte superficial del rotacional es
$\nn\times(\vP_2-\vP_1)\,\delta_S$: es nula si y solo si se conserva la componente tangencial de $\vP$. Con esto
podemos responder a la pregunta de la \cref{obs:D-no-determinado}.

\begin{proposicion}[Cuándo se anula $\vD$]\label{prop:D-nulo}
Si en todo el espacio no hay carga libre y la polarización es irrotacional, incluidas sus superficies de
discontinuidad,
\[
  \rhof\equiv0,\qquad \rot\vP=\vcero\ \text{en cada medio},\qquad \nn\times(\vP_2-\vP_1)=\vcero\ \text{en cada interfase},
\]
entonces $\vD\equiv\vcero$ y, por tanto, $\vE=-\vP/\eo$. Recíprocamente, si $\vD\equiv\vcero$, no hay carga libre y $\vP$
es irrotacional.
\end{proposicion}

\begin{proof}
Por \eqref{eq:gauss-D} y \eqref{eq:rot-D}, $\vD$ tiene divergencia y rotacional nulos en todo el espacio, incluidos
los términos superficiales. Si el dieléctrico es acotado, fuera de él $\vD=\eo\vE$ decae como el campo de una
distribución neutra (\cref{prop:QP-nula}), al menos como $1/r^3$. Por el \cref{lem:helmholtz}, el único campo con
estas propiedades es el nulo. El recíproco es inmediato: si $\vD=\vcero$, $\rhof=\divg\vD=0$ y $\rot\vP=\rot\vD=\vcero$.
\end{proof}

\begin{ejemplo}[Polarizaciones radiales]\label{ej:radial}
Una polarización $\vP=f(r)\,\uv r$ (esférica) o $\vP=f(\rho)\,\uv\rho$ (cilíndrica) es irrotacional en el interior del
material, y en sus superficies, que son esferas o cilindros coaxiales, $\vP$ es normal y $\nn\times\vP=\vcero$. Sin
carga libre, $\vD=\vcero$ en todo el espacio y $\vE=-\vP/\eo$ dentro del material: así se resuelven los
problemas~\ref{prob:electrete-esferico} y~\ref{prob:electrete-coaxial}.
\end{ejemplo}

\begin{ejemplo}[El cilindro uniformemente polarizado]\label{ej:stokes}
En un cilindro con $\vP$ uniforme y axial, $\rot\vP=\vcero$ dentro, pero en la pared lateral $\vP$ es tangente y salta a
cero. Un pequeño rectángulo que atraviese la pared (\cref{fig:p44}\,(a)) tiene circulación $\oint\vP\cdot\dl=P\,\Delta\ell\neq0$ y,
por el teorema de Stokes, $\rot\vP\neq\vcero$ en la pared. Aunque no haya carga libre, $\vD\neq\vcero$
(\cref{prob:cilindro-polarizado}): sus líneas se cierran como las del campo de un solenoide
(\cref{fig:cilindro-lineas}).
\end{ejemplo}

\begin{trampa}[Trampa habitual]
«Como no hay carga libre, $\vD=\vcero$» solo es cierto si además $\vP$ es irrotacional (\cref{prop:D-nulo}). En el
cilindro, la esfera o el cubo uniformemente polarizados, $\vD\neq\vcero$ aunque $\rhof=0$.
\end{trampa}
